\subsection{Origin of the Einstein tensor}
  \label{subsec:einstein-origin}

  The crucial structural point is the following.

  The coefficient $a_2(x)=R/6$ implies that the pole subtraction generates a
  local term proportional to
  \begin{equation}
    \int d^4x\,\sqrt{g}\, R .
  \end{equation}

  Upon metric variation, this counterterm produces the Einstein tensor,
  \begin{equation}
    \frac{\delta}{\delta g^{\mu\nu}}
    \int d^4x\,\sqrt{g}\, R
    \;\propto\;
    G_{\mu\nu} .
  \end{equation}

  Therefore, the local two-derivative Einstein term in the renormalized
  spectral stress tensor arises from the counterterm associated with the
  $a_2$ heat-kernel pole.

  It does not arise from the bare finite part of $\zeta'_{A_g}(0)$ alone.
  Rather, it is generated through renormalization of the geometric sector,
  in the standard sense of induced gravity.

  This mechanism is structurally identical to Sakharov's induced-gravity
  construction~\cite{Sakharov1967}, with the spectral cutoff scale playing the role of the
  ultraviolet regulator.

  Higher-derivative local contributions originate from $a_4$ and yield
  four-derivative tensors (Bach-like structures) upon variation.
  Non-local form factors arise from the finite part of the determinant
  and will be treated separately.
